% !PRINTED 578
\weberpage{578}
%<<CONT
Zahlenkoeffizienten rational sind. Daß es ganze Zahlen sind, ist
wegen des Nenners 2, der zunächst auftritt, nicht zu ersehen.

Um dies näher zu untersuchen, betrachten wir die Wurzeln
der Funktion $A(x)$: Sie sind, wenn gesetzt wird:
\[
\Omega_{h,h'} = \frac{2hK+2h'iK'}{m},
\]
\weq{9}{x_{h,h'} = \sqrt{\varkappa}\,\mathrm{sn}\,\Omega_{h,h'} = \frac{\vartheta_{11}\left(\frac{h+h'\omega}{m},\omega\right)}{\vartheta_{01}\left(\frac{h+h'\omega}{m},\omega\right)},}
worin $h,h'$ irgend welche ganze Zahlen sein können; nur dürfen
sie nicht beide gleich Null sein. Man erhält \gesperrt{alle Wurzeln} von
$A(x)$, wenn man $h,h'$ je ein volles Restsystem nach dem Modul
$m$ durchlaufen läßt, abgesehen von der Kombination 0,0. Es ist
dann
\weq{10}{x_{h,h'} = \frac{-i\sum\limits_{-\infty,\infty}^{\nu}(-1)^\nu e^{\frac{\pi i h}{m}(2\nu+1)} q^{\left(\frac{2h+1}{2}+\frac{h'}{m}\right)^2}}{\sum\limits_{-\infty,\infty}^{\nu}(-1)^\nu e^{\frac{2\pi i \nu h}{m}} q^{\left(\nu+\frac{h'}{m}\right)^2}}.}
% [source sic: the exponent prints 2h+1 although the summation index suggests 2\nu+1]
Um diesen Ausdruck nach steigenden Potenzen von $q$ zu entwickeln,
muß man zunächst die niedrigste Potenz von $q$ im Nenner
aufsuchen, d. h. das Glied, in dem $\nu+\frac{h'}{m}$ so klein als möglich
wird. Dies findet statt, wenn $\nu$ die zunächst an $-h'/m$ gelegene
ganze Zahl ist. Es gibt, da $m$ ungerade ist, \gesperrt{nur eine} solche
Zahl $\nu$, und die Formel (10) erhält die Gestalt:
\weq{11}{x_{h,h'} = \frac{Q_1}{1+Q_2} = Q_1 - Q_1 Q_2 + Q_1 Q_2^2 - \dots,}
worin $Q_1$ und $Q_2$ nach steigenden Potenzen von $q$ geordnete Reihen
sind, deren Koeffizienten ganze algebraische Zahlen (Kreisteilungszahlen)
sind.

\gesperrt{Demnach läßt sich $x_{h,h'}$ nach steigenden Potenzen
von $q$ entwickeln und die Koeffizienten dieser Entwickelung
sind ganze algebraische Zahlen.}

Ist $S$ eine symmetrische Funktion der $x_{h,h'}$ mit rationalen
ganzzahligen Koeffizienten, so läßt auch diese sich in derselben
Weise nach Potenzen von $q$ entwickeln. Andererseits ist $S$ nach
dem Fundamentalsatz über symmetrische Funktionen rational durch
%CONT>>
